% ===========================================================================
\section{Compiler correctness: $\VD\iso\IR(\VD)$}\label{sec:compiler}
% ===========================================================================

\subsection{The translation}

\begin{definition}[Compiler]\label{def:compiler}
Let $A=(\Loc,\ell_0,\Clk,\Ch,\Edges,\Inv)$ satisfy \cref{rem:side}, with
$\Loc=(\ell_0',\dots,\ell_{k-1}')$, $\Clk=(x_1,\dots,x_n)$ and $\Edges=(e_1,\dots,e_m)$. Define
the index maps
\[
  \nu_\Loc(\ell_i')=i,\qquad \nu_\Clk(x_i)=i,\ \nu_\Clk(0)=0,\qquad \nu_\Edges(e_j)=t_{j-1},
\]
and lift $\nu_\Clk$ to atoms and conjunctions:
$\nu(x-y\bowtie c)=(\nu_\Clk(x),\nu_\Clk(y),\bowtie,c)$ and $\nu(G)=\{\nu(\varphi)\mid\varphi\in G\}$.
Let $\mathit{sort}$ denote the canonical list of a set. Then $C(A)=M$ with
\begin{itemize}
\item $X=(x_1,\dots,x_n)$, $\Ch_M=\mathit{sort}(\Ch)$, $R$ and $\mathit{info}$ from the options;
\item $\lambda_i=(\mathrm{name}(\ell_i'),\ \mathit{sort}(\nu(\Inv(\ell_i'))))$;
\item $\iota=\nu_\Loc(\ell_0)$ $(=0$ by \cref{rem:side});
\item for $e_j=(\ell,g,\alpha,Y,\ell')$:
  $t_{j-1}=(\mathit{tid}_{j-1},\nu_\Loc(\ell),\nu_\Loc(\ell'),\alpha,\mathit{sort}(\nu(g)),\mathit{sort}(\nu_\Clk(Y)))$,
  where $\mathit{tid}$ is ``$\mathit{src}.\alpha.\mathit{tgt}$'', suffixed by $\#r$ for the
  $r$-th repetition of the same triple in document order;
\item $p_i=(\at{\mathrm{name}(\ell_i')},i,I_D(\ell_i'))$, with $I_D(\ell)$ the entry of the
  DT interpretation file or the empty string.
\end{itemize}
This is \code{twin::compiler::compile\_file} (\code{src/compiler/compiler.cpp}).
\end{definition}

\begin{lemma}\label{lem:nu}
$\nu_\Loc$, $\nu_\Clk$ and $\nu_\Edges$ are bijections onto the locations, clock
identifiers and transitions of $C(A)$, and $C(A)$ is well-formed.
\end{lemma}
\begin{proof}
One IR location, clock and transition is emitted per source location, clock
and edge, in the same order. Location and clock names are unique in $A$
(UTAP rejects duplicates, and the compiler re-checks), so the inverse maps are
well defined. Transition identifiers are unique because repetitions are
numbered. Well-formedness is checked by \code{ir::validate} before the IR is
emitted. A failure there would be a compiler defect, reported as diagnostic
\code{TWC090}, and no IR would be produced.
\end{proof}

\begin{lemma}[Constraint preservation]\label{lem:constraint}
For every conjunction $G$ over $\Clk$ and valuation $v$:
$v\sat G \iff (v\circ\nu_\Clk^{-1})\sat\mathit{sort}(\nu(G))$.
\end{lemma}
\begin{proof}
Let $w=v\circ\nu_\Clk^{-1}$, so $w(\nu_\Clk(x))=v(x)$, and $w(0)=v(0)=0$. For an atom,
$w(\nu_\Clk(x))-w(\nu_\Clk(y))=v(x)-v(y)$, so $v\sat\varphi\iff w\sat\nu(\varphi)$.
Satisfaction of a conjunction depends only on the \emph{set} of its atoms, and
$\mathit{sort}$ does not change that set.
\end{proof}

\subsection{The isomorphism}

\begin{theorem}[Compiler correctness]\label{thm:compiler}
Let $A$ satisfy \cref{rem:side}, be initialisable, and let $M=C(A)$. The map
\[
  \Gamma:\Valid(A)\to\Valid(M),\qquad
  \Gamma(\ell,v,\theta)=(\nu_\Loc(\ell),\ v\circ\nu_\Clk^{-1},\ \theta)
\]
is a bijection, and $(\Gamma,\nu_\Edges,\pi)$ with $\pi(\at{\ell})=p_{\nu_\Loc(\ell)}$ is an
isomorphism $\sem{A}\iso\sem{M}$ (\cref{def:iso}). The isomorphism maps internal
labels to internal labels and preserves action labels:
$\mathit{act}(e)=\alpha(\nu_\Edges(e))$.
\end{theorem}
\begin{proof}
\emph{Bijection.} $\Gamma$ is the product of bijections (\cref{lem:nu}) with
identity on $\theta$. The side conditions $v(x)\le\theta\le\Hz$ are
invariant under renaming. Validity is preserved and reflected by
\cref{lem:constraint} applied to $\Inv(\ell)$ and $J_{\nu_\Loc(\ell)}=\mathit{sort}(\nu(\Inv(\ell)))$.

\emph{(i) Initial state.} $\Gamma(\ell_0,\mathbf{0},0)=(\nu_\Loc(\ell_0),\mathbf{0},0)=(\iota,\mathbf{0},0)$.
$A$ is initialisable iff $M$ is, by \cref{lem:constraint}.

\emph{(ii) Delays.} By \cref{lem:convex}, $(\ell,v,\theta)\trans{d}(\ell,v+d,\theta+d)$
holds iff $\theta+d\le\Hz$, $v\sat\Inv(\ell)$ and $v+d\sat\Inv(\ell)$. By
\cref{lem:constraint}, and since $(v+d)\circ\nu_\Clk^{-1}=(v\circ\nu_\Clk^{-1})+d$, this
is exactly the premise of \textsc{(Delay$_M$)} for $\Gamma(\ell,v,\theta)$. The targets
correspond: $\Gamma(\ell,v+d,\theta+d)=(\nu_\Loc(\ell),(v\circ\nu_\Clk^{-1})+d,\theta+d)$.

\emph{(iii) Discrete steps.} Let $e=(\ell,g,\alpha,Y,\ell')$ and $t=\nu_\Edges(e)$.
The premises of \textsc{(Disc)} are $v\sat g$ and $v[Y{:=}0]\sat\Inv(\ell')$. Those of
\textsc{(Disc$_M$)} at $\Gamma(\ell,v,\theta)$ are $\nu_\Loc(\ell)=\mathit{src}(t)$, which holds
by construction, $w\sat\mathit{sort}(\nu(g))$, and $w[\nu_\Clk(Y){:=}0]\sat J_{\nu_\Loc(\ell')}$.
They coincide by \cref{lem:constraint}, using
$v[Y{:=}0]\circ\nu_\Clk^{-1}=(v\circ\nu_\Clk^{-1})[\nu_\Clk(Y){:=}0]$. The targets
correspond in the same way. Conversely, every $t\in T$ is $\nu_\Edges(e)$ for a
unique $e$, so every IR step is the image of a source step.

\emph{(iv) Propositions.} $\lab_A(\ell,v,\theta)=\{\at{\ell}\}$ and
$\lab_M(\Gamma(\ell,v,\theta))=\{p_{\nu_\Loc(\ell)}\}=\{\pi(\at{\ell})\}$.

Actions are copied verbatim, so internal edges map to internal transitions.
\end{proof}

\subsection{From the abstract compiler to the implemented one}

\Cref{thm:compiler} is about the mathematical function $C$. Two links connect
it to the bytes on disk.

\begin{assumption}[Source reading]\label{ass:utap}
UTAP parses the UPPAAL document into the abstract syntax of
\cref{def:source}. This is the same assumption the alignment verdict rests on,
because the aligner uses the same parser.
\end{assumption}

\begin{proposition}[Translation validation]\label{prop:tv}
Suppose compilation succeeds. Then the automaton $A_c$ read by the compiler and
the automaton $A_a$ read by the aligner's \code{dtpta::TimedAutomaton} agree on:
the template, the clock sequence and numbering, the location names and their
order, the initial location (index $0$), the edge sequence (endpoints, actions,
reset sets), and, for every location and edge, the \emph{set of valuations}
satisfying the invariant and the guard. Consequently $\sem{A_c}=\sem{A_a}$.
\end{proposition}
\begin{proof}
These are exactly the checks of \code{crosscheck\_with\_aligner}
(\code{src/compiler/crosscheck.cpp}). Compilation fails with \code{TWC060} if any
of them fails. Invariants and guards are compared as zones: both readings are
intersected with the universe of non-negative valuations, closed, and compared
with UDBM's \code{dbm\_areEqual}. Two closed, non-empty DBMs are equal iff they
denote the same zone; empty results are compared as empty. The semantics of
\cref{def:source-sem} depends on conjunctions only through their satisfying
valuations, and all other components are equal, so the two LTTSs coincide.
\end{proof}

\begin{corollary}\label{cor:compiled-is-verified}
If compilation succeeds, then $\sem{\IR(\VD)}\iso\sem{A_a}$, where $A_a$ is the
automaton the aligner analysed (up to the aligner's own zone semantics; see
\cref{sec:limitations}).
\end{corollary}
\begin{proof}
\Cref{thm:compiler} gives $\sem{A_c}\iso\sem{M}$ and \cref{prop:tv} gives
$\sem{A_c}=\sem{A_a}$.
\end{proof}

\begin{remark}[Determinism of compilation]
$C$ is a function of the source bytes, the interpretation file and the options
$(\text{model id},\text{version},R)$. The implementation uses no timestamps,
paths, environment variables or iteration over unordered containers in the IR
(the compilation \emph{manifest} contains a timestamp and is not hashed into the
IR). Recompiling therefore yields a byte-identical IR. The test suite checks
this on the aligner's corpus.
\end{remark}
